\section{Formal connections and the two ideal identities}
\label{sec:connections}

We retain the notation of Section~\ref{sec:deformation}.  All
connections and integrals in this section have characteristic-zero
formal coefficients.  We construct polarized Chern--Simons functions
on the walls and prove two ideal identities.  Their internal derivatives generate
the vertex defect ideal, and the total wall value on the restriction
graph belongs to the square of that ideal. The first conclusion
transfers the cubic holonomy laws to gradient-ideal memberships; the
second permits the graph correction in Section~\ref{sec:algebra}.
After proving these identities, we control all finite gradient jets
in one complex algebraic ring.  This is the input for algebraic
replacement; the full Chern--Simons series need not be algebraic.

\subsection{Realizing transition cochains by connections}

Off the triangle equations, the edge matrices need not be transition
functions of a bundle.  We first replace them by a genuine smooth
formal bundle and connection while preserving the exact flat family
modulo the triangle-defect ideal.  The explicit inverses in the
construction will also control its finite algebraic jets.

\begin{lemma}[The image bundle]
\label{conn:image-bundle}
For a finite algebraic transition-cochain chart on a compact smooth
manifold \(Z\), let \(\mathfrak d\) be its triangle-defect ideal.
There is a smooth formal rank-two bundle and connection whose
reduction modulo \(\mathfrak d\) is the specified flat family.
Its fiber curvature lies in the closed extension of \(\mathfrak d\).
The construction permits a total connection including parameter
directions.  If the reference bundle is trivial, it has a global
formal frame extending any chosen reference frame.
\end{lemma}

\begin{proof}
Choose smooth functions \(\phi_\alpha\) subordinate to the good
cover, with support contained in their patches and
\(\sum_\alpha\phi_\alpha^2=1\).  In a sufficiently large trivial
bundle form the block matrix
\[
 e_{\alpha\beta}=\phi_\alpha\phi_\beta g_{\alpha\beta}.
\]
Terms with disjoint supports are zero; support containment makes
the extensions by zero smooth.  The triangle identities imply
\(e^2-e\in\mathfrak d\).  Put \(\Delta=e^2-e\) and define
\begin{equation}
 E=\frac{1+(2e-1)(1+4\Delta)^{-1/2}}2 .
 \label{conn:idempotent}
\end{equation}
The binomial series is defined formally because
\(\Delta\in\mathfrak m\); indeed our charts have
\(\mathfrak d\subset\mathfrak m^2\).
All factors commute with \(e\), and
\((2e-1)^2=1+4\Delta\), so \(E^2=E\).
Moreover \(E=e\) modulo \(\mathfrak d\).
Its image has rank two, as its reference image does.

On a patch indexed by \(\alpha\), set
\[
 C_\alpha=(\phi_\beta g_{\beta\alpha})_\beta,\qquad
 R_\alpha=(\phi_\beta g_{\alpha\beta})_\beta,\qquad
 S_\alpha=EC_\alpha .
\]
The first is a block column, the second a block row.  One has the
exact identity
\begin{equation}
 R_\alpha C_\alpha
       =\sum_\beta\phi_\beta^2g_{\alpha\beta}g_{\beta\alpha}
       =1.
 \label{conn:row-column}
\end{equation}
Also \(eC_\alpha=C_\alpha\) modulo \(\mathfrak d\), so
\[
 T_\alpha:=R_\alpha S_\alpha-1\in\mathfrak d,\qquad
 L_\alpha=(1+T_\alpha)^{-1}R_\alpha\big|_{\operatorname{im}E}.
\]
These give a frame and its inverse on the rank-two image bundle.
For example \(L_\alpha S_\alpha=1\), and a rank-two frame with
this left inverse is an isomorphism onto the rank-two image.
In the reduction, the frames satisfy
\(S_\beta=S_\alpha g_{\alpha\beta}\) on intersections.

Choose a subordinate partition \((\chi_\alpha)\) and blend the
connections which are trivial in these frames:
\begin{equation}
 \nabla s=\sum_\alpha\chi_\alpha
                    S_\alpha\,d_{\mathrm{tot}}(L_\alpha s).
 \label{conn:blended-connection}
\end{equation}
This is a connection since each summand satisfies the Leibniz
rule and \(\sum\chi_\alpha=1\).  In the reduction modulo
\(\mathfrak d\) its fiber part is exactly the flat connection
defined by the constant-in-fiber transitions.  Its fiber
curvature therefore has coefficients in \(\mathfrak d\).
Using the full differential in \eqref{conn:blended-connection}
also supplies a total connection, regardless of the parameter
dependence of the reduced transition matrices.

If a reference global frame is given, view it as a matrix of
sections of \(\operatorname{im}E(0)\), extend its coefficients
constantly in the parameters, and apply \(E\).  The resulting
sections reduce to that frame, and hence are a formal global
frame.  Inverses exist by the maximal-ideal Neumann expansion.
\end{proof}

We shall also use a basic lifting property implicit in this
construction.  A formal flat connection over a complete local
quotient has parallel frames on a contractible patch, uniquely
specified by their values at one point.  To verify this, pass
to every finite parameter quotient, solve the ordinary parallel
transport equations with values in its finite-dimensional
coefficient algebra, and use flatness and contractibility for
path independence.  These constructions commute with quotient
maps and give compatible formal frames.  Thus a flat bundle
isomorphism of reductions can be tested by based parallel
transport, including over nonreduced coefficient rings.

\begin{lemma}[Prescribed wall collars]
\label{conn:wall-collar}
For a wall chart \((a_i,y_i)\), the connection in
Lemma~\ref{conn:image-bundle} can be chosen, in a formal frame on
a torus collar, to have total connection form
\begin{equation}
 U_i(a_i)\,d\theta_i+V_i(a_i,y_i)\,d\eta_i ,
 \label{conn:collar}
\end{equation}
with zero parameter components.  It still reduces to the given
flat family modulo \(\mathfrak d_i\).  The collar frame extends
to a global formal frame on \(J_i\) with \(A_i(0)=0\).
\end{lemma}

\begin{proof}
Temporarily write \(U,V,\mathfrak d\) for the wall data.
Modulo \(\mathfrak d\) the two torus holonomies commute.
Their formal logarithms commute as well: they are power series
in the two commuting matrices.  On the simply connected cover
of the collar, take the flat parallel frame \(P\) issuing from
the prescribed based frame.  In the chosen convention the
frame
\[
 P\exp(\theta U+\eta V)
\]
is periodic, since the exponential cancels the deck
holonomies \(\exp(-U)\), \(\exp(-V)\).  Its fiber connection
form is \(U\,d\theta+V\,d\eta\).

To lift it, use finitely many small collar patches with lifted
torus coordinates.  On each patch a lift of \(P\) is expressed
using an image frame \(S_\alpha\) and a fixed finite transition
word.  Multiply by the displayed exponential.  These local
lifted frames have the same reduction modulo \(\mathfrak d\);
their partition-of-unity average is therefore a formal frame
with that reduction.  In this frame prescribe
\eqref{conn:collar} near the boundary, including its zero
parameter components, and use a collar cutoff to interpolate
with the previous total connection.  Their fiber reductions
are the same flat connection.  Thus the interpolation preserves
the reduction and keeps the fiber curvature in \(\mathfrak d\).
This reasoning does not require the off-equation matrices
\(U,V\) to commute.

At the origin the representation on each wall component is
trivial, and the based choices make the collar frame the
restriction of a global parallel frame.  Lift that frame
globally as in Lemma~\ref{conn:image-bundle}.  On a collar its
discrepancy from the prescribed frame is a matrix
\(G=1+O(\mathfrak m)\).  Multiplication of the global frame by
\(1+\chi(G-1)\), extending coefficients by zero using a collar
cutoff \(\chi\), makes it agree with the prescribed frame on a
shorter collar.  This matrix is formally invertible.  The
resulting global frame has the asserted properties.  Closed
wall components require only the initial global parallel
frame and its formal lift.
\end{proof}

\subsection{Polarized Chern--Simons functions}

The functional below uses the classical Chern--Simons form
\cite[Section~3, Propositions~3.2 and~3.8]{ChernSimons1974}
with the normalization displayed here.
The boundary polarization and the identities needed in this proof
are derived below.

\begin{definition}[The polarized wall function]
\label{conn:potential-definition}
Use the global frame in Lemma~\ref{conn:wall-collar}, and denote
the fiber connection form by \(A_i\).  With
\(\epsilon_i=\int_{\partial J_i}d\theta_i\wedge d\eta_i\), put
\begin{align}
 \operatorname{CS}_{J_i}(A_i)
   &=\int_{J_i}\operatorname{tr}
        \left(A_i\wedge dA_i+\frac23 A_i\wedge A_i\wedge A_i\right),
       \notag\\
 P_i(a_i,y_i)
   &=\operatorname{CS}_{J_i}(A_i)
                  +\epsilon_i\operatorname{tr}(U_iV_i).
 \label{conn:polarized-potential}
\end{align}
The differential in this integral is fiberwise.  Integrals over
all closed components are included with their chosen orientations.
\end{definition}

\begin{lemma}[Polarized variation]
\label{conn:variation}
Writing \(F_A=dA+A\wedge A\), every parameter variation satisfies
\begin{equation}
 \delta P_i
   =2\int_{J_i}\operatorname{tr}(\delta A_i\wedge F_{A_i})
          +2\epsilon_i\operatorname{tr}(V_i\,\delta U_i).
 \label{conn:variation-formula}
\end{equation}
Consequently \(P_i\in\mathfrak m_i^3\) and every coefficient of
its parameter differential lies in
\(\mathfrak k_i\mathbb C[[a_i,y_i]]\).
\end{lemma}

\begin{proof}
Differentiate the integrand in
\eqref{conn:polarized-potential}.  Cyclic trace, with the signs
from degrees of forms, gives
\[
 \delta\operatorname{tr}
      \left(A\wedge dA+\frac23 A^3\right)
  =2\operatorname{tr}(\delta A\wedge F_A)
                           -d\operatorname{tr}(A\wedge\delta A).
\]
On the torus the integrated boundary term is
\[
 -\epsilon_i\operatorname{tr}
                 (U_i\,\delta V_i-V_i\,\delta U_i).
\]
Adding the variation of
\(\epsilon_i\operatorname{tr}(U_iV_i)\) gives
\eqref{conn:variation-formula}.
The fiber curvature is in \(\mathfrak d_i\subset\mathfrak m_i^2\),
and \(V_i\in\mathfrak m_i^2\).  Thus each first derivative of
\(P_i\) has order at least two.  Since \(A_i(0)=0\), also
\(P_i(0)=0\), which proves \(P_i\in\mathfrak m_i^3\).
Finally \(\mathfrak d_i\) and \(V_i\) are contained in the
completed ideal \(\mathfrak k_i\); integration preserves that
ideal by Definition~\ref{def:coefficient-convention}.
\end{proof}

\subsection{The derivative ideal}

We now compare the wall functions on the two boundary copies of every
cut. Their internal derivatives will recover every generator of the
vertex defect ideal, including its nonreduced structure.  The reason is
cohomological: internal wall variations detect the restricted curvature
classes by duality, and the degree-two Mayer--Vietoris isomorphism
detects the vertex curvature class.  We apply this to each possible
first obstruction to lifting the flat vertex family.

Define the separated function and its restriction by
\begin{equation}
 \Psi(a,C)=\sum_{i=1}^d
       \bigl(P_i(a_i,C_i^+)-P_i(0,C_i^-)\bigr),
 \qquad C=(C_i^+,C_i^-)_{i=1}^d\in\mathcal Y ,
 \label{conn:psi}
\end{equation}
and use \(Y=(Y_i^+,Y_i^-)_{i=1}^d\) from
Proposition~\ref{def:restriction-charts}.  A notation such as
\(\partial_C\Psi(a,Y)\) always means: first differentiate with
respect to the independent \(C\)-coordinates, then substitute
\(C=Y(a,t)\).

\begin{lemma}[A traceless vertex curvature]
\label{conn:vertex-connection}
The vertex chart admits a global formal connection \(A_X\)
whose curvature is traceless and belongs to \(\mathfrak r\).
It reduces to the exact transition family modulo
\(\mathfrak r\).  Its reference bundle is trivial.
\end{lemma}

\begin{proof}
Apply Lemma~\ref{conn:image-bundle}.  The reference bundle
comes from a graph; a complex rank-two bundle with connected
structure group on a graph is trivial, by choosing a frame
on its vertices and extending over its edges.  Its pullback
to each component of \(X\) is therefore trivial.  A continuous
reference frame may be smoothly approximated in the smooth
bundle, preserving invertibility, and then lifted formally.

Modulo \(\mathfrak r\), the special-linear transitions give
a parallel determinant section.  Lift it to a nonvanishing
formal section of the determinant bundle; local coefficient
lifts and a partition suffice, and the reference section
is nonzero.  If its induced determinant connection has
one-form \(\alpha\), then \(\alpha\in\mathfrak r\), because
the reduced section is parallel.  Subtract
\(\frac12\alpha\,1\) from the rank-two connection.
This preserves its reduction and makes the determinant
section parallel, so the resulting curvature has trace zero.
\end{proof}

\begin{proposition}[Generation of the full defect ideal]
\label{conn:derivative-ideal}
In \(B=\mathbb C[[a,t]]\),
\begin{equation}
 \bigl(\partial_C\Psi(a,Y(a,t))\bigr)=\mathfrak r.
 \label{conn:derivative-identity}
\end{equation}
\end{proposition}

\begin{proof}
In a \(y_i\)-variation, \(U_i\) is fixed.  Thus the boundary
term of \eqref{conn:variation-formula} is zero.
After either substitution \(\iota_i^\pm\), the wall curvature
belongs to \(\mathfrak r\), since the restriction is exactly
flat modulo \(\mathfrak r\).  All the derivatives in
\eqref{conn:derivative-identity} consequently belong to
\(\mathfrak r\).

It remains to prove that their images span
\(\mathfrak r/\mathfrak m\mathfrak r\).
Suppose a nonzero functional \(\lambda\) on this space
annihilates all those images.  Form its small extension
\eqref{def:small-extension}, so that
\(I=\mathfrak r/\mathfrak s\) is one-dimensional and
\(\mathfrak mI=I^2=0\).
Reduce the connection from Lemma~\ref{conn:vertex-connection}
to \(B/\mathfrak s\).  Its curvature is an \(I\)-valued
traceless two-form.  Since multiplying \(I\) by a
positive-order parameter coefficient gives zero, Bianchi
becomes the reference equation
\(d_{A_X(0)}F_{A_X}=0\).  The curvature thus defines
\[
 \kappa\in H^2(X;\mathfrak l_\rho)\otimes_{\mathbb C} I.
\]

The wall connections and the restrictions of \(A_X\)
are isomorphic modulo \(\mathfrak r\), by the based
restriction construction.  Lift a smooth formal
identification of their bundles over \(B/\mathfrak s\).
In that identification their difference is an
\(I\)-valued one-form \(\eta\).  Curvature changes by
exactly \(d_{A(0)}\eta\): its quadratic term is zero since
\(I^2=0\), and the positive-order part of \(A\) acts
trivially since \(\mathfrak mI=0\).
Thus the wall curvature classes are the restrictions
of \(\kappa\), in the common reference frames.

At the wall origin a \(y_i\)-derivative
\(\partial_{y_{i,j}}A_i(0)\) is a closed one-form, because
wall curvature has order at least two.  Its tangential
boundary value is zero: \(U_i\) does not vary with \(y_i\),
and \(V_i\) has zero linear term.  Its absolute cohomology
class is the corresponding chart tangent, as is seen by
first-order holonomy.  Hence these forms represent all
of \(H^1(J_i,\partial J_i;\mathfrak l)\), by
\eqref{def:relative-identification}.
If a general-linear frame adds a scalar exact term,
traceless projection has the same \(\mathfrak l\)-class
and the same pairing with the traceless curvature.
The assertion includes every closed wall component.

Now apply \(\lambda\) to the internal derivatives.
The polarized variation formula reduces to the
nonzero scalar \(2\) times the trace integration pairing
of these relative one-forms with the wall restrictions
of \(\kappa\), with an irrelevant overall minus sign
on negative copies.  Poincare--Lefschetz duality
\cite[Theorem~3.43]{Hatcher2002} gives
a perfect pairing of these relative classes with
\(H^2(J_i;\mathfrak l)\): the manifolds are compact
oriented three-manifolds and the trace form identifies
the coefficient system with its dual.  Thus every wall
restriction of \(\kappa\) vanishes.
Apply the degree-two isomorphism of
Lemma~\ref{def:paired-mv} to the pair \((\kappa,0)\).
Its image is zero, so \(\kappa=0\) on every component
of \(X\).

De Rham cohomology with a flat smooth coefficient bundle
computes this local-system cohomology.  Consequently
there is a traceless reference one-form \(\eta_X\),
with values in \(I\), satisfying
\[
 d_{A_X(0)}\eta_X=-F_{A_X}.
\]
Then \(A_X+\eta_X\) is \emph{exactly} flat over
\(B/\mathfrak s\); all omitted nonlinear terms are zero
by \(\mathfrak mI=I^2=0\).  It has the prescribed
reduction over \(B/\mathfrak r\).

On every contractible patch choose the original parallel
frame of that reduction, namely the one realizing its
chart transitions.  Lift its initial value at a point
and parallel-transport with the corrected connection.
Uniqueness of parallel transport makes the reduced frame
equal to the original specified frame.  The lifted
constant transition matrices therefore give a flat lift
of the actual vertex cochain, not merely of an unbased
isomorphism class.  This contradicts
Lemma~\ref{def:equation-chart}.  The contradiction does
not require the corrected connection to remain in the
chosen \(U\)-slice; the equation-chart lemma excludes
arbitrary cochain corrections.  No nonzero annihilating
functional exists.

The quotient \(\mathfrak r/\mathfrak m\mathfrak r\)
is finite-dimensional.  Its derivative images therefore
span, and Nakayama's lemma applied to the quotient of
\(\mathfrak r\) by the derivative ideal proves
\eqref{conn:derivative-identity}.
\end{proof}

\subsection{The square of the ideal and gluing}

The derivative identity is now established. To control the value of
the potential on the restriction graph, we compare connections on the
closed boundary of $X$. Changes of connection and changes of frame
have different effects on Chern--Simons values; the next lemma records
both before they are used in the gluing argument.

\begin{lemma}[Two Chern--Simons comparisons on a closed manifold]
\label{conn:closed-comparisons}
Let \(N\) be a closed oriented three-manifold, and let \(A\)
be a global formal connection with \(F_A\in I\), for a
scalar ideal \(I\).
\begin{enumerate}
\item If \(\eta\in I\) is a one-form, then
\[
 \operatorname{CS}_N(A+\eta)-\operatorname{CS}_N(A)\in I^2.
\]
\item For an arbitrary smooth formal gauge \(g\),
\[
 \operatorname{CS}_N(A^g)-\operatorname{CS}_N(A)
\]
is independent of the parameters.  The reference gauge
need not be nullhomotopic.
\end{enumerate}
\end{lemma}

\begin{proof}
For the first assertion put \(A_u=A+u\eta\),
\(0\leq u\leq1\).  Its curvature
\(F_A+u\,d_A\eta+u^2\eta\wedge\eta\) lies in \(I\).
The variation formula on the closed manifold gives
\[
 \operatorname{CS}_N(A+\eta)-\operatorname{CS}_N(A)
     =2\int_0^1\int_N\operatorname{tr}(\eta\wedge F_{A_u})\,du
                                                        \in I^2.
\]
For the second use
\(A^g=g^{-1}Ag+g^{-1}d_Ng\).
For a parameter variation set \(\xi=g^{-1}\delta g\).
Then
\[
 \delta A^g=g^{-1}(\delta A)g+d_{A^g}\xi,\qquad
 F_{A^g}=g^{-1}F_Ag.
\]
Bianchi and the closedness of \(N\) give
\[
 \delta\bigl(\operatorname{CS}_N(A^g)
                         -\operatorname{CS}_N(A)\bigr)
     =2\int_N d_N\operatorname{tr}(\xi F_{A^g})=0.
\]
In characteristic zero a formal series with zero first
derivatives is constant.  This argument permits any
reference homotopy class of \(g\); its possible degree
affects only that constant.
\end{proof}

\begin{proposition}[The square-ideal value identity]
\label{conn:value-ideal}
For the same maps and potential,
\begin{equation}
 \Psi(a,Y(a,t))\in\mathfrak r^2 .
 \label{conn:value-identity}
\end{equation}
\end{proposition}

\begin{proof}
All wall connections in this proof are first pulled back
by their maps \(\iota_i^\pm\) to \(B=\mathbb C[[a,t]]\).
Their curvature lies in \(\mathfrak r\).
On a positive wall the second torus slope bounds in
\(Q\), so \(V_i(a_i,Y_i^+)\in\mathfrak r\).
Choose a fiber cutoff \(\chi\) supported on its collar,
equal to one near the boundary, and put
\[
 A_{i,u}=A_i-u\chi V_i\,d\eta_i,\qquad 0\leq u\leq1 .
\]
The connection change lies in \(\mathfrak r\), and
every interpolated curvature also lies in
\(\mathfrak r\).  This is immediate from the curvature
change formula and does not assume commutation away
from the defect quotient.
Use the polarized functional along this path:
\[
 P_{i,u}=\operatorname{CS}_{J_i}(A_{i,u})
                  +\epsilon_i\operatorname{tr}(U_i(1-u)V_i).
\]
Here \(U_i\) is held fixed.  Formula
\eqref{conn:variation-formula} has zero boundary term
for the \(u\)-variation; its bulk term is a product of
two \(\mathfrak r\)-valued factors.  Thus
\(P_{i,1}-P_{i,0}\in\mathfrak r^2\).
At \(u=1\) the polarization correction is zero and
the collar form is the pure form \(U_i\,d\theta_i\).
On a negative wall \(U_i(0)=0\) already, so its
functional is raw Chern--Simons and its collar form
is the pure form \(V_i(0,Y_i^-)\,d\eta_i\).
It follows that, modulo \(\mathfrak r^2\),
\(\Psi(a,Y)\) is the oriented sum of the raw wall
Chern--Simons integrals of these modified connections.

We next construct an \emph{exactly flat} formal
connection on \(Q\) with those pure collar forms.
Its reduction must be the actual restriction from
the flat vertex family, including based transports.
Fix the standard paths from the main basepoint of
\(Q\) to the seam basepoints.  Lift their comparison
transport matrices from \(B/\mathfrak r\) to \(B\).
Assign each free end generator the required pure
collar holonomy, conjugated by its lifted comparison
transport and adjusted for its orientation.
The other slope at that seam bounds in \(Q\).
Lift the matrices of the \(h\)-generators from the
vertex restriction arbitrarily in
\(\operatorname{SL}_2(B)\).  Such lifts exist because
the group is smooth at every reference matrix.
Since the end generators and \(h\)-generators freely
generate \(\pi_1(Q)\), these choices impose no
relations and define a representation over \(B\)
lifting the specified one over \(B/\mathfrak r\).
Its associated flat bundle has an exactly flat
connection.  Choose its frame at each seam basepoint
using the lifted comparison transport; parallel
transport followed by a one-slope exponential gives
the prescribed pure connection form on that collar.

These collar frames extend to a global frame on
\(Q\).  At the origin this is a relative section
problem for a special-linear bundle on a compact
three-manifold, with prescribed section on its
boundary collars and at its basepoint.
The group \(\operatorname{SL}_2(\mathbb C)\) retracts
onto \(\operatorname{SU}_2\cong S^3\), so its
\(\pi_0,\pi_1,\pi_2\) vanish.  The obstruction to
extending a section across a relative cell of
dimension \(j\leq3\) lies in \(\pi_{j-1}\) of
this group; all such obstructions vanish.
A smooth reference extension follows by relative
smooth approximation.  At successive formal orders
the discrepancy is identity modulo the maximal
ideal and extends coefficientwise with a cutoff,
just as in Lemma~\ref{conn:wall-collar}.
This yields a global formal frame agreeing with
every prescribed collar frame.

In this frame \(\operatorname{CS}_Q\) is constant in
the parameters.  Its variation has zero bulk term
by exact flatness.  On each torus both its connection
form and its variation lie along the same single
one-form \(d\theta_i\) or \(d\eta_i\), so
\(\operatorname{tr}(A_Q\wedge\delta A_Q)\) is zero.
The boundary term of the raw Chern--Simons variation
therefore also vanishes.

Glue these framed pieces to form a connection
\(A_{\partial}\) in a global frame on the closed
three-manifold \(\partial X\).  The forms agree on
the product collars, so their integrals add.
We must compare the assembled flat reduction with
the actual restriction from \(X\), including its
transports.  The component \(X_0=X_Q\) containing \(Q\)
contains all the torus-bearing wall copies in its
boundary.  Each such copy attaches to \(Q\) along
exactly one connected torus.  The incidence graph
has one \(Q\)-vertex and one leaf for each such
wall copy.  It has no cycles.
Fix the flat comparison on \(Q\), including its
seam-basepoint values.  On a wall its based
representation is the prescribed restriction,
so choose the flat comparison with that same
value at its single seam basepoint.
Parallel transport makes the two comparisons
agree on the entire connected torus.
There is no second attachment imposing an
additional transport condition.  Every closed
wall component is an isolated boundary piece
for this purpose and is compared separately,
whether it belongs to \(X_0\) or to another
component of \(X\).  Thus the assembled reduction
is isomorphic to the full boundary restriction
of the vertex family.

Lift this boundary bundle isomorphism formally.
In global frames this amounts to lifting smooth
matrix coefficients; the reference determinant
is nonzero, so the lifted matrix has a formal
inverse.  No extension of this identification
over \(X\) is required.  Denote by \(g\) the
resulting boundary gauge comparison with the
already chosen global frame of \(A_X\).
Then in a common boundary frame
\[
 A_{\partial}=A_X^g|_{\partial X}+\eta,
                         \qquad \eta\in\mathfrak r .
\]
Both compared curvatures are in \(\mathfrak r\).
The first part of Lemma~\ref{conn:closed-comparisons}
shows that this \emph{connection change} contributes
only an element of \(\mathfrak r^2\).
The second part shows that the separate
\emph{frame change} contributes a parameter-constant
term, which disappears after subtracting reference
values.  The gauge \(g(0)\) may have nonzero degree;
it is not asserted to extend over \(X\).

Finally apply Stokes in the global frame already
chosen on \(X\).  The identity
\[
 d\,\operatorname{tr}
      \left(A_X\wedge dA_X+\frac23 A_X^3\right)
                 =\operatorname{tr}(F_{A_X}\wedge F_{A_X})
\]
follows by the graded cyclic trace rule.  Hence
\[
 \operatorname{CS}_{\partial X}(A_X|_{\partial X})
       =\int_X\operatorname{tr}(F_{A_X}\wedge F_{A_X})
                                                  \in\mathfrak r^2.
\]
Stokes applies componentwise to the compact
oriented smooth four-manifold with its rounded
collared corners, including any components with
empty boundary.  The \(Q\)-contribution is
constant, all \(P_i(0)\) are zero, and the preceding
connection and gauge comparisons hold after
subtracting the origin values.  These facts prove
\eqref{conn:value-identity}.
\end{proof}

\subsection{Algebraic control of the wall gradients}

The two ideal identities are now proved in the completed parameter
ring.  To pass to algebraic replacement, we return to the original
wall functions \(P_i\) and the total connections that define them.
The modified connections in the square-ideal proof were comparison
devices after pullback to the vertex chart; they did not change these
wall functions.  We show that every finite gradient jet has a
representative in one fixed smooth complex algebraic ring.

\begin{proposition}[All gradient jets in one complex algebraic ring]
\label{conn:all-jets}
For each wall there is one fixed algebraic chart ring \(A_i^*\),
with completion \(\mathbb C[[a_i,y_i]]\), containing the original
transition germs and \(\mathfrak k_i\), such that for every
integer \(b\geq1\), every coefficient of \(d_{\mathrm{param}}P_i\)
modulo \(\mathfrak k_i^b\) is the image of an element of \(A_i^*\).
The ring is independent of \(b\).
\end{proposition}

\begin{proof}
Fix a wall \(Z=J_i\), put \(v=(a_i,y_i)\), and write
\(B=\mathbb C[[v]]\), \(\mathfrak d=\mathfrak d_i\), and
\(\mathfrak k=\mathfrak k_i\).  Choose the common etale chart
ring \(A^*\) once, containing the finite transition cochain, its
defects, and the two holonomy matrices \(H_{i,1},H_{i,2}\).

The global frame defining \(P_i\) has only formal control in the
maximal ideal; this does not give algebraic coefficients modulo
powers of \(\mathfrak k\).  We first express the gradient without
derivatives of that frame, then control the local expressions
that remain.

Suppress the wall index.  In the global frame defining \(P\), write
the total connection as
\[
 \nabla=d_Z+d_{\mathrm{param}}+A+\Gamma .
\]
Here \(A\) and \(\Gamma\) are its fiber and parameter components.
For a parameter vector \(\delta\), define the mixed curvature
with the parameter argument first:
\begin{equation}
 K_\delta=F_{\mathrm{tot}}(\delta,\,\cdot\,)
                    =\delta A-d_A\Gamma_\delta .
 \label{conn:mixed-curvature}
\end{equation}
Bianchi's identity \(d_AF_A=0\) and trace invariance give
\[
 \operatorname{tr}(d_A\Gamma_\delta\wedge F_A)
                      =d_Z\operatorname{tr}(\Gamma_\delta F_A).
\]
Consequently the exact polarized variation is
\begin{align}
 \delta P
  &=2\int_Z\operatorname{tr}(K_\delta\wedge F_A)
    +2\int_{\partial Z}\operatorname{tr}(\Gamma_\delta F_A)
    +2\epsilon\operatorname{tr}(V\,\delta U) \notag\\
  &=2\int_Z\operatorname{tr}(K_\delta\wedge F_A)
                          +2\epsilon\operatorname{tr}(V\,\delta U).
 \label{conn:controlled-gradient}
\end{align}
The second equality holds because the global frame agrees
with the prescribed collar frame for every parameter, and
the total connection has \(\Gamma_\delta=0\) there.  Both
curvature components in the bulk trace transform by
conjugation under every parameter-dependent frame change.
We can therefore compute that trace in the local image and
collar frames of Lemmas~\ref{conn:image-bundle} and
\ref{conn:wall-collar}, without differentiating the global gauge.

We now check that the coefficients needed in
\eqref{conn:controlled-gradient} can be computed in \(A^*\) to
any prescribed \(\mathfrak k\)-adic order.  On each smooth
coordinate patch we maintain the following property:
modulo \(\mathfrak k^b\), each required coefficient is a finite
sum
\begin{equation}
 \sum_{\nu=1}^{N_b} f_\nu(x)\,a_\nu(v),
       \qquad f_\nu\text{ smooth},\quad a_\nu\in A^* .
 \label{conn:finite-sums}
\end{equation}
The functions, constants, and number of summands may depend
on \(b\); the complex algebra \(A^*\) does not.
We verify this for every operation actually used.

Recall the image-frame construction:
\(S_\alpha=EC_\alpha\),
\(T_\alpha=R_\alpha S_\alpha-1\in\mathfrak d\), and
\(L_\alpha=(1+T_\alpha)^{-1}R_\alpha|_{\operatorname{im}E}\).
The block matrix \(e_{\alpha\beta}=\phi_\alpha\phi_\beta
g_{\alpha\beta}\) has the finite-sums property exactly.  Since
\(\Delta=e^2-e\in\mathfrak d\subset\mathfrak k\),
the binomial series in \eqref{conn:idempotent} truncates after
finitely many terms modulo \(\mathfrak k^b\).  The explicit
inverse \eqref{conn:row-column} is equally important.  In the
notation of Lemma~\ref{conn:image-bundle},
\[
 (1+T_\alpha)^{-1}
       \equiv\sum_{j=0}^{b-1}(-T_\alpha)^j
                              \pmod{\mathfrak k^b}.
\]
Thus image frames and their inverses satisfy
\eqref{conn:finite-sums}; no smooth partition function has
been inverted.  In frame \(S_\beta\), the total connection
\eqref{conn:blended-connection} has matrix
\[
 \sum_\alpha\chi_\alpha
        (L_\beta S_\alpha)\,
                   d_{\mathrm{tot}}(L_\alpha S_\beta).
\]
This expression introduces only the already specified inverse
matrices, multiplication, and differentiation.

Next consider the averaged collar frame.  Its local lifts have
the form
\[
 T_\gamma=S_{\alpha(\gamma)}W_\gamma
                       \exp(\theta_\gamma U+\eta_\gamma V),
\]
where \(W_\gamma\) is a fixed transition word.  Both \(U,V\)
belong to \(\mathfrak kB\), since the corresponding matrix
deviations generate part of \(\mathfrak k\).  The logarithms
and this exponential consequently truncate to finite
expressions modulo each \(\mathfrak k^b\).  Write
\(T=\sum_\gamma\chi_\gamma T_\gamma\) for the average.
On a fixed small patch choose one participating lift \(T_0\).
All the lifts agree modulo \(\mathfrak d\); hence
\[
 T-T_0=\sum_\gamma\chi_\gamma(T_\gamma-T_0)
                                           \in\mathfrak d.
\]
Modulo \(\mathfrak k\), the exponential is \(1\).
An explicit leading inverse on that patch is therefore
\[
 V_0=W_0^{-1}R_{\alpha(0)}
                       \big|_{\operatorname{im}E}.
\]
It satisfies \(V_0T=1+N\) with \(N\in\mathfrak k\), and the
inverse of the average is
\begin{equation}
 T^{-1}\equiv
       \left(\sum_{j=0}^{b-1}(-N)^j\right)V_0
                                      \pmod{\mathfrak k^b}.
 \label{conn:average-inverse}
\end{equation}
This proves the finite-sums property for the average and its
inverse.  Agreement of its reductions, not merely invertibility
of its summands, is what permits this calculation.
The collar interpolation uses only fixed smooth cutoffs and
these expressions, so it preserves the same property.

The coordinate derivations of the affine parameter space
extend uniquely to its etale algebra and preserve its local
ring \(A^*\).  In an implicit presentation $\mathcal F(v,w)=0$, differentiation
gives
\[
 \partial_{v_j}w
 =-\bigl(\partial_w\mathcal F(v,w)\bigr)^{-1}
                           \partial_{v_j}\mathcal F(v,w).
\]
The variable Jacobian determinant is a unit germ in \(A^*\), because
its value at the origin is nonzero.  Its inverse therefore already
belongs to this local etale ring.
The product rule gives
\begin{equation}
 \partial_{v_j}(\mathfrak k^{b+1}B)
                                     \subset\mathfrak k^bB.
 \label{conn:derivative-order}
\end{equation}
Fiber derivatives preserve scalar ideal order.  Accordingly,
compute primitive frame and logarithm expressions one ideal
order higher when taking a parameter derivative.

The fiber connection and fiber curvature use only fiber
derivatives of primitive frame expressions.
The parameter connection component uses one parameter
derivative, and \(K_\delta\) uses at most one: its two terms
are a parameter derivative of \(A\) and a fiber covariant
derivative of \(\Gamma_\delta\).
Thus calculations modulo \(\mathfrak k^{b+1}\), followed by
\eqref{conn:derivative-order}, suffice to compute every term
of \eqref{conn:controlled-gradient} modulo \(\mathfrak k^b\).
The torus term is likewise controlled by finite logarithm
expansions one order higher.

Finally take a finite partition to integrate the local
expressions.  The integral of \eqref{conn:finite-sums} is
\(\sum_\nu(\int f_\nu)a_\nu\), which belongs to \(A^*\)
because every integration constant is a complex scalar.
For the remainder, work modulo \(\mathfrak m^N\).
It has values in the finite-dimensional scalar subspace
which is the image of \(\mathfrak k^b\), and integration
preserves that subspace.  The integrated remainder therefore
lies in \(\mathfrak k^b+\mathfrak m^N\) for every \(N\).
Ideals of a complete Noetherian local ring are adically
closed, so it lies in \(\mathfrak k^b\).
This justifies the assertion with the completed smooth
coefficient modules actually in use.
\end{proof}

\begin{remark}[Scope of the fixed-ring assertion]
\label{conn:fixed-ring-scope}
Proposition~\ref{conn:all-jets} asserts finite-sum jets in one
local \emph{complex} algebraic ring.  It does not assert that an
arbitrary integral of algebraic functions of a parameter is
algebraic, or that all coefficients of the original
Chern--Simons series lie in one finitely generated
\(\mathbb Z\)-algebra.  The latter kind of coefficient ring
will be chosen only after algebraic replacement and a finite
selection of ideal-membership witnesses.
\end{remark}

\subsection{The data supplied to algebraic replacement}

We collect the identities just proved and transfer the cube-zero
holonomy law to the two types of functions used by the final algebraic
obstruction.  Both the first-slope coordinates and the negative-wall
first-slope derivatives lie in the end-deviation ideal modulo
\(\mathfrak r\).

\begin{corollary}[The deformation output]
\label{conn:output}
The algebraic maps in
Proposition~\ref{def:restriction-charts} and the formal functions
\(P_i\in\mathbb C[[a_i,y_i]]\) have the following properties:
\begin{enumerate}
\item \(P_i\in\mathfrak m_i^3\), its gradient is zero modulo
      \(\mathfrak k_i\), and all its gradient jets have
      representatives in one fixed complex algebraic
      chart ring as in Proposition~\ref{conn:all-jets}.
\item The two tangent images of \(Y(0,t)\) and
      \(\sigma Y(0,t)\) are complementary \(D\)-planes
      in the \(2D\)-dimensional \(C\)-space.
\item For \(\Psi\) in \eqref{conn:psi},
\[
  \bigl(\partial_C\Psi(a,Y)\bigr)=\mathfrak r,\qquad
                   \Psi(a,Y)\in\mathfrak r^2 .
\]
\item Choose one scalar first-slope coordinate \(a_{i,1}\)
      for each \(i\), and define
\begin{equation}
 p_i(y_i)=\partial_{a_{i,1}}P_i(0,y_i).
 \label{conn:p-functions}
\end{equation}
      For every triple of labels, including repetitions,
\begin{equation}
 a_{i,1}a_{j,1}a_{k,1},\qquad
 p_i(Y_i^-)\,p_j(Y_j^-)\,p_k(Y_k^-)
                 \ \in\ \bigl(\partial_C\Psi(a,Y)\bigr).
 \label{conn:cubic-laws}
\end{equation}
\end{enumerate}
Moreover the cube of the sum of the pulled-back wall
ideals is contained in \(\mathfrak r\), as in
Proposition~\ref{def:cubic-nilpotence}.
\end{corollary}

\begin{proof}
Only the last item remains to be checked.  Work modulo
\(\mathfrak r\), where the end-deviation ideal
\(\mathfrak n\) satisfies \(\mathfrak n^3=0\).
Each \(a_{i,1}\) is an algebraic coordinate of an end
holonomy at the identity, and hence belongs to
\(\mathfrak n\).  Evaluate
\eqref{conn:variation-formula} with the
\(a_{i,1}\)-variation on the negative copy,
at \((0,Y_i^-(a,t))\).  Its curvature term vanishes
modulo \(\mathfrak r\), while its torus term is
\[
 2\epsilon_i\operatorname{tr}
       \left(V_i(0,Y_i^-)\,
                    \partial_{a_{i,1}}U_i(0)\right).
\]
The remaining second-slope matrix on that copy
is the end holonomy, with the prescribed orientation
and basing.  Its negative logarithm consequently
has entries in \(\mathfrak n\), so
\(p_i(Y_i^-)\in\mathfrak n\).
Both kinds of triple products vanish modulo
\(\mathfrak r\).  Proposition~\ref{conn:derivative-ideal}
identifies that ideal with the derivative ideal,
which gives \eqref{conn:cubic-laws}.
\end{proof}
